%%% SECTION 3: CALIFORNIA AND FEDERAL REGULATORY LANDSCAPE %%%

California serves as the site of the first Physical AI oncology clinical trial,
making a dual-jurisdiction analysis of state and federal law essential for
understanding the complete regulatory environment. This section adapts the
comprehensive California and federal legal framework analysis to Physical AI
oncology trial contexts, providing the regulatory map that trial sponsors,
investigators, and sites must navigate.

\subsection{Regulatory Classification Framework}
\label{subsec:regulatory-classification}

The regulatory landscape for Physical AI in oncology trials can be organized
along two classification axes. The first axis distinguishes between AI systems
that qualify as regulated medical products and those that function as research
operational tools without direct medical product classification. The second axis,
applicable to device-classified Physical AI systems, distinguishes between
investigational and approved status, and between significant risk and
non-significant risk designations.

The federal device definition under 21 U.S.C. \S~321(h) encompasses instruments,
machines, and other articles intended for use in diagnosis, treatment, or
prevention of disease, which includes many Physical AI system components. The
investigational device exemption authority under 21 U.S.C. \S~360j(g)
establishes procedures for device use during clinical investigation. The FD\&C
Act \S~520(o) software exclusions from the 21st Century Cures Act exclude
certain software functions from the device definition, but these exclusions were
designed primarily for software-only clinical decision support and may not
encompass the full range of Physical AI functionalities that involve direct
physical patient interaction.

Physical AI systems present classification challenges because a single system
may simultaneously occupy multiple categories. A surgical robot performing a
tumor resection functions as a medical device. The AI algorithm guiding its
surgical decisions may qualify as software as a medical device (SaMD). The data
collection and transmission functions may be classified as electronic data
systems under Part 11. The robot's interaction with the patient triggers human
subjects protection requirements. This multi-category nature underscores the
need for the integrated regulatory approach provided by this National Platform.

\subsection{California Regulatory Authority}
\label{subsec:california-authority}

\subsubsection{Human Subjects Protections}

California's Protection of Human Subjects in Medical Experimentation Act (Cal.
Health \& Safety Code \S\S~24170-24179.5) establishes protections for research
participants that supplement federal requirements. The Experimental Subject's
Bill of Rights (\S~24172) enumerates specific rights that must be communicated
to research participants, creating obligations beyond federal informed consent
requirements. The definition of ``medical experiment'' (\S~24174) is broad
enough to encompass Physical AI procedures where robots perform functions that
constitute medical experimentation.

For Physical AI trial participants in California, these provisions create
additional protections. The Experimental Subject's Bill of Rights requires
disclosure of the nature of the experiment (including robotic involvement),
risks and discomforts (including robot malfunction scenarios), and the right to
withdraw. The Clinical Trial Site Complete Documentation \cite{site-docs}
incorporates these California-specific requirements into its documentation
package, and AB 2847 from that package establishes additional patient rights
specific to Physical AI interactions.

\subsubsection{Health Data Privacy and Security}

The Confidentiality of Medical Information Act (CMIA, Cal. Civ. Code
\S\S~56-56.37) governs the handling of medical information by healthcare
providers and other entities. Electronic medical information safeguards
(\S~56.101) require specific protections for electronically maintained medical
data, directly applicable to the digital records generated by Physical AI
systems. California's data breach notification requirements (Cal. Civ. Code
\S~1798.82) and medical information breach requirements (Cal. Health \& Safety
Code \S~1280.15) impose notification obligations when Physical AI system
security breaches compromise patient data.

Physical AI systems generate new categories of health data including robot
telemetry, sensor readings, video recordings of procedures, and AI decision
logs. CMIA requirements apply to these data categories when they constitute
medical information, creating compliance obligations that interact with federal
HIPAA \cite{hipaa} protections. The Physical AI Federated Learning framework
\cite{fl-paper} addresses these dual privacy requirements through
privacy-preserving approaches that minimize data centralization.

\subsubsection{Consumer Privacy}

CCPA/CPRA (Cal. Civ. Code \S\S~1798.100 et seq.) provides broad consumer
privacy protections with specific exemptions for clinical trial and biomedical
research data. Clinical trial data may be statutorily exempt in many common
trial contexts, but non-exempt data streams from Physical AI systems may still
trigger CCPA obligations. For example, if a Physical AI system collects data
about a patient's movements or preferences that falls outside the clinical trial
data exemption, CCPA rights to access, deletion, and opt-out would apply.

\subsubsection{AI-Specific Healthcare Communications}

AB 3030 (Cal. Health \& Safety Code \S~1339.75) requires disclosure when
AI-generated content is used in patient clinical communications. AB 489
addresses health advice from artificial intelligence. Physical AI systems that
communicate with patients, whether through social companion robots providing
emotional support or through AI-generated treatment explanations, must comply
with these disclosure requirements. The patient instructions documented in
Patient Instructions: Physical AI Trials \cite{patient-instructions-paper}
incorporate appropriate disclosures for each of the ten robot types.

\subsubsection{Medical Practice and Licensing}

California's unlicensed practice of medicine provisions (Cal. Bus. \& Prof.
Code \S~2052) raise questions about Physical AI systems that perform functions
traditionally reserved for licensed practitioners. When a surgical robot
performs tissue manipulation or a needle-placement robot conducts a biopsy, the
boundary between machine-assisted medical practice and autonomous machine
practice becomes legally significant. The adapted 21 CFR Part 312
\cite{cfr312-adapt} addresses this through its Physical AI operator definition,
which establishes that qualified human oversight must be maintained even when
robots perform physical procedures.

\subsection{Federal Regulatory Authority}
\label{subsec:federal-authority}

\subsubsection{Investigational Product Frameworks}

The IND regulations (21 CFR Part 312) \cite{cfr-part-312} govern investigational
drug use in clinical trials. When Physical AI systems are used alongside
investigational drugs, the IND framework must accommodate the additional
complexity of robotic drug administration, AI-guided dosing, and automated
safety monitoring. The Physical AI Adaption: 21 CFR Part 312 \cite{cfr312-adapt}
provides the adapted framework spanning ten subparts, including the Phase 0
simulation validation stage that requires demonstration of Physical AI
system safety before human subjects enrollment.

IDE regulations (21 CFR Part 812) govern investigational device use. When
Physical AI systems themselves qualify as investigational devices, sponsors must
obtain IDE approval or qualify for exemption. The interaction between IND and
IDE requirements becomes particularly complex when Physical AI systems are used
to administer investigational drugs, potentially triggering both regulatory
pathways simultaneously.

\subsubsection{Human Subjects Protection}

Federal human subjects protection operates through 21 CFR Part 50 \cite{cfr-part-50}
(FDA-specific protections), 21 CFR Part 56 (IRB requirements), and the Common
Rule (45 CFR Part 46). Physical AI adds new dimensions to informed consent
because participants must understand their interactions with robotic systems, AI
decision-making processes, safety monitoring systems, and their right to refuse
or pause robotic procedures. The Physical AI Adaption: 21 CFR Part 50
\cite{cfr50-adapt} addresses these dimensions through 17 Physical AI-specific
definitions and a new Subpart C dedicated to Physical AI safety requirements.

The NCI Central Institutional Review Board \cite{nci-cirb} provides a model for
centralized IRB review that could be adapted for Physical AI trial oversight,
enabling consistent evaluation of Physical AI risks and protections across
multiple trial sites.

\subsubsection{Electronic Records and Data Integrity}

21 CFR Part 11 governs electronic records and electronic signatures in
FDA-regulated activities. Physical AI systems generate, store, and transmit
electronic clinical trial records including robot telemetry, sensor data,
autonomous decision logs, imaging data, and treatment records. These constitute
new categories of electronic records that must comply with Part 11 requirements
for trustworthiness and reliability.

Interoperability standards including FHIR \cite{hl7-fhir} for healthcare data
exchange and DICOM \cite{dicom} for medical imaging provide standardized
formats for some Physical AI data streams. The MCP server infrastructure
\cite{national-mcp-paper} extends these standards by providing standardized
communication protocols for Physical AI-specific data that existing standards
do not cover, including real-time robot telemetry, task order management, and
safety interlock states.

\subsubsection{Privacy and Security}

HIPAA Privacy and Security Rules (45 CFR Parts 160 and 164) \cite{hipaa} govern
the use and disclosure of protected health information (PHI) by covered
entities and business associates. Physical AI systems handling PHI must comply
with HIPAA requirements, creating obligations for data encryption, access
controls, audit trails, and breach notification. The federated learning
framework \cite{fl-paper} and federated learning in healthcare
\cite{federated-learning-healthcare} provide privacy-preserving approaches
that reduce the need to centralize PHI while enabling multi-site model
improvement.

\subsection{Robotics-Specific Regulatory Considerations}
\label{subsec:robotics-regulatory}

FDA's position on robotically-assisted surgical (RAS) systems in oncology
presents a specific regulatory consideration. RAS systems have not received
marketing authorization specifically for cancer prevention or treatment,
although they are used in oncology-related procedures under existing
clearances. The significance of oncology-related endpoints (recurrence,
disease-free survival, overall survival) for Physical AI trials means that
trial designs must demonstrate not just procedural safety but also oncology
outcome equivalence or superiority.

The Physical AI Adaption: 21 CFR Part 312 \cite{cfr312-adapt} addresses this
through its Phase 0 simulation validation requirement, which allows Physical
AI systems to demonstrate safety and preliminary efficacy through simulation
before any real patient contact. The 24-hour simulation evidence from the
Clinical Trial Site Complete Documentation \cite{site-docs} provides a model
for how such Phase 0 validation could be conducted at scale.

\subsection{Gaps and Opportunities}
\label{subsec:gaps-opportunities}

Several regulatory gaps and ambiguities exist for Physical AI implementations
that this National Platform addresses. First, the lack of specific FDA guidance
for Physical AI clinical trials means that sponsors must navigate multiple
existing guidances and regulations without a unified reference. This platform
provides that unified reference through its comprehensive body of adapted
standards and documentation.

Second, jurisdictional overlaps between California and federal requirements
create compliance complexity, particularly for data privacy where CMIA, CCPA,
and HIPAA all may apply to the same data. The site establishment documentation
\cite{site-docs} addresses this by integrating both jurisdictions into a single
compliance framework.

Third, classification challenges arise for multi-function Physical AI systems
that simultaneously qualify as medical devices, electronic data systems, and
research tools. The adapted regulatory standards provide classification
guidance, with the 21 CFR Part 312 adaptation \cite{cfr312-adapt} establishing
which Physical AI functions require IND coverage and which fall under device
or operational tool categories.

Fourth, workforce transition considerations must be addressed as Physical AI
enables treating more patients with more robots and fewer workers. California
labor law, federal employment regulations, and healthcare staffing requirements
all intersect with the workforce implications of Physical AI adoption. The
implementation strategy in Section~\ref{sec:implementation} addresses these
workforce considerations through phased deployment with training and role
evolution programs.

\clearpage

\begin{table}[H]
\centering
\caption{Regulatory Gaps and National Platform Solutions}
\label{tab:gaps-solutions}
\small
\begin{tabularx}{\textwidth}{X X}
\toprule
\textbf{Regulatory Gap} & \textbf{National Platform Solution} \\
\midrule
No unified Physical AI trial guidance & Adapted standards (ICH, CFR 50, CFR 312) \\
Dual-jurisdiction compliance complexity & Integrated site documentation package \\
Multi-function classification challenges & Physical AI definitions across all adaptations \\
Workforce transition regulatory uncertainty & Phased implementation with training programs \\
Absence of Physical AI safety standards & PSL (3 dimensions) and USL (4 dimensions) \\
No multi-site Physical AI infrastructure & MCP servers and federated learning \\
\bottomrule
\end{tabularx}
\end{table}

The regulatory landscape for Physical AI oncology trials is complex but
navigable. By building on existing frameworks rather than creating entirely new
ones, the adapted standards in this platform \cite{ich-adapt, cfr50-adapt,
cfr312-adapt} provide a credible and practical path forward that the
pharmaceutical and regulatory industries can adopt with confidence.

\subsection{Comparative Obligation Analysis}
\label{subsec:comparative-obligations}

The dual-jurisdiction challenge requires sponsors and sites to satisfy both
California state and federal requirements simultaneously. Table~\ref{tab:comparative-obligations}
maps the key regulatory domains to their California and federal counterparts,
identifying where requirements overlap, complement, or create unique obligations
for Physical AI trials.

\begin{longtable}{
    >{\raggedright\arraybackslash}p{2.95cm} 
    >{\raggedright\arraybackslash}p{4.8cm} 
    >{\raggedright\arraybackslash}p{4.5cm} 
    >{\raggedright\arraybackslash}p{2.95cm}
}
\caption{California vs. Federal Regulatory Obligations for Physical AI Trials}
\label{tab:comparative-obligations} \\ 
\toprule 
\textbf{Domain} & \textbf{California Requirement} & \textbf{Federal Requirement} & \textbf{Interaction} \\
\midrule
\endfirsthead
\toprule
\textbf{Domain} & \textbf{California Requirement} & \textbf{Federal Requirement} & \textbf{Interaction} \\
\midrule
\endhead
\bottomrule
\endfoot 
Human subjects protection & Cal. Health \& Safety Code \S\S~24170-24179.5 (Experimental Subject's Bill of Rights) & 21 CFR Part 50 (informed consent); 45 CFR Part 46 (Common Rule) & California adds rights beyond federal minimum \\ 
Health data privacy & CMIA (Cal. Civ. Code \S\S~56-56.37); electronic safeguards (\S~56.101) & HIPAA Privacy and Security Rules (45 CFR Parts 160, 164) & Both apply; CMIA may be stricter for certain data categories \\
Consumer data privacy & CCPA/CPRA (Cal. Civ. Code \S\S~1798.100 et seq.) & No federal equivalent for non-health consumer data & CCPA applies to non-exempt Physical AI data streams \\
AI communications & AB 3030 (AI-generated patient communications disclosure); AB 489 (health advice from AI) & No federal equivalent & California-specific disclosure obligations for Physical AI patient interactions \\
Medical practice & Cal. Bus. \& Prof. Code \S~2052 (unlicensed practice) & FD\&C Act device/drug definitions and authorized uses & Physical AI operator definition bridges both frameworks \\
Electronic records & State data breach notification (Cal. Civ. Code \S~1798.82) & 21 CFR Part 11 (electronic records/signatures) & Both apply to Physical AI system records \\
Device regulation & State medical device requirements & 21 CFR Parts 312 (IND), 812 (IDE), device clearance & Federal preemption for device regulation; state adds protections \\
IRB review & California research review requirements & 21 CFR Part 56; 45 CFR Part 46 & Both required; NCI CIRB model may streamline \\
Trial registration & California clinical trial transparency laws & 42 CFR Part 11 (ClinicalTrials.gov) & Both require registration; federal database is primary \\
\end{longtable}

The comparative analysis reveals that California's regulatory environment
provides additional protections beyond federal requirements in several key
areas. The Experimental Subject's Bill of Rights (\S~24172) creates California-
specific patient rights that supplement federal informed consent. CMIA
provisions for electronic medical information (\S~56.101) add state-level
security requirements for Physical AI-generated health data that complement
HIPAA. The AI-specific communications laws (AB 3030 and AB 489) have no
federal equivalents, meaning Physical AI systems that communicate with patients
must comply with these California-specific disclosure requirements in addition
to any federal guidance.

For trial sponsors planning Physical AI trials in California, the dual-
jurisdiction framework requires compliance with the stricter of the two
requirements in each domain. The site establishment documentation \cite{site-docs}
integrates both jurisdictions into a unified compliance checklist, and the
adapted 21 CFR Part 50 \cite{cfr50-adapt} incorporates informed consent
elements that satisfy both California and federal requirements. When Physical
AI trials expand to other states (Phase 2 and Phase 3 of the implementation
strategy in Section~\ref{sec:implementation}), each state's regulatory
landscape must be similarly analyzed and integrated into the compliance
framework.

\subsection{FDA Clearance Pathways for Physical AI in Oncology}
\label{subsec:fda-clearance-pathways}

Physical AI systems may require regulatory clearance through multiple FDA
pathways depending on their classification and intended use. Surgical robots
used in oncology trials may be cleared through 510(k) premarket notification
(if a suitable predicate device exists), de novo classification (for novel
device types), or premarket approval (PMA) for Class III devices. The adapted
21 CFR Part 312 \cite{cfr312-adapt} accommodates all these pathways through its
Physical AI System Description requirement in the IND submission, which
documents the regulatory status of each Physical AI system.

The regulatory pathway distinction is practically significant: a da Vinci
surgical robot with existing FDA clearance for general surgical use may be
deployed in a Physical AI oncology trial under its current clearance, with the
trial protocol specifying oncology-specific use parameters. By contrast, a
novel Physical AI system without prior clearance may require IDE approval before
deployment in clinical trials. The Phase 0 simulation validation requirement
provides additional safety evidence that may support both IND submissions (for
drug trials using Physical AI) and IDE applications (for Physical AI device
investigations), creating regulatory efficiency across both pathways.

%%% ADDITIONAL CONTENT FOR SECTION 3 %%%

\subsubsection{Expanded Comparative Obligation Analysis by Regulatory Domain}

The comparative obligation framework established in
Table~\ref{tab:comparative-obligations} merits further elaboration across
several domains where the interaction between California and federal
requirements creates particularly complex compliance obligations for Physical
AI trial sponsors. This expanded analysis addresses the practical
implementation considerations that sponsors must resolve when designing
Physical AI compliance systems.

In the data privacy domain, the interaction between CMIA, CCPA/CPRA, and
HIPAA creates a layered compliance obligation with distinct requirements at
each layer. HIPAA establishes the baseline through its Privacy Rule (45 CFR
Part 164 Subpart E), which governs the use and disclosure of protected health
information by covered entities and business associates. CMIA adds
California-specific requirements that exceed HIPAA in several respects: CMIA
covers medical information from a broader range of entities than HIPAA's
covered entity definition, CMIA's electronic safeguard requirements
(\S~56.101) impose specific technical obligations beyond HIPAA's
``reasonable safeguard'' standard, and CMIA's breach notification timeline
(15 business days) is more aggressive than HIPAA's 60-day notification
window. CCPA/CPRA adds consumer privacy rights that apply to non-exempt
Physical AI data streams, including data about patient preferences,
scheduling patterns, and non-clinical interactions with Physical AI systems.

For Physical AI trials, the practical implication is that sponsors must
implement data governance systems capable of satisfying all three frameworks
simultaneously. Robot telemetry data that constitutes medical information
must comply with both CMIA and HIPAA. AI decision logs that contain protected
health information must meet HIPAA security standards while also satisfying
CMIA electronic safeguard requirements. Non-clinical data collected by
Physical AI systems (such as facility navigation patterns or environmental
sensor readings that do not constitute medical information) may still be
subject to CCPA/CPRA if the data can be linked to identifiable consumers.
The federated learning framework \cite{fl-paper} reduces the complexity of
this layered compliance by minimizing data centralization, thereby reducing
the volume of data subject to each framework's requirements.

Table~\ref{tab:privacy-detailed} provides a detailed comparison of specific
privacy obligations across the three frameworks as they apply to Physical AI
data categories.

\clearpage

\begin{table}[H]
\centering
\caption{Detailed Privacy Obligations for Physical AI Data Categories}
\label{tab:privacy-detailed}
\small
\begin{tabularx}{\textwidth}{l c c c}
\toprule
\textbf{Physical AI Data Category} & \textbf{HIPAA} & \textbf{CMIA} & \textbf{CCPA/CPRA} \\
\midrule
Patient vital signs (robot-captured) & PHI & Medical info & Exempt \\
Robot telemetry (position, force) & PHI if linked & Medical info if clinical & May apply \\
Surgical video recordings & PHI & Medical info & Exempt \\
AI diagnostic outputs & PHI & Medical info & Exempt \\
AI treatment recommendations & PHI & Medical info & Exempt \\
Autonomous decision logs & PHI if linked & Medical info if clinical & May apply \\
Patient scheduling data & PHI if linked & Not medical info & Applies \\
Facility sensor readings & Not PHI & Not medical info & May apply \\
Equipment maintenance logs & Not PHI & Not medical info & Not applicable \\
De-identified research data & Not PHI & Not medical info & Not applicable \\
Federated model parameters & Not PHI & Not medical info & Not applicable \\
\bottomrule
\end{tabularx}
\end{table}

The classification of Physical AI data categories in
Table~\ref{tab:privacy-detailed} reveals that the privacy framework
applicable to a given data element depends on both the nature of the data
and its linkage to identifiable individuals. Robot telemetry data
(position, force, speed) is not inherently protected health information, but
becomes PHI under HIPAA and medical information under CMIA when linked to a
specific patient during a clinical procedure. This means that the same data
stream may transition between privacy categories depending on context,
requiring Physical AI data systems to implement context-aware access controls
and audit mechanisms. The MCP data flow server \cite{national-mcp-paper}
addresses this through its metadata tagging system, which associates privacy
classification markers with each data element at the point of capture.

In the informed consent domain, the expanded analysis reveals that Physical
AI trials in California must address at least five distinct consent
obligations simultaneously. First, 21 CFR 50.25 requires eight basic
elements and six additional elements of informed consent for FDA-regulated
research. Second, the Experimental Subject's Bill of Rights (Cal. Health \&
Safety Code \S~24172) requires specific disclosures including the nature of
the experiment, reasonably foreseeable risks, expected benefits, available
alternative procedures, and the right to withdraw. Third, AB 3030 requires
disclosure when AI-generated content is used in patient clinical
communications, directly applicable to Physical AI systems that generate
treatment explanations or status updates. Fourth, AB 489 addresses health
advice from artificial intelligence, applicable when Physical AI systems
provide health-related guidance to trial participants. Fifth, the adapted
21 CFR Part 50 \cite{cfr50-adapt} establishes 17 Physical AI-specific
definitions and requires disclosure of robotic system specifications, AI
decision-making processes, safety monitoring systems, and emergency stop
procedures.

The cumulative consent disclosure requirements create a tension between
comprehensiveness and comprehensibility. Research has consistently
demonstrated that overly long and detailed consent documents reduce patient
understanding and may undermine the informed consent process. The patient
instructions framework \cite{patient-instructions-paper} addresses this
tension by providing layered information: a concise consent document that
meets all legal requirements is supplemented by detailed, robot-specific
patient instructions that participants can review at their own pace. This
layered approach satisfies both the legal obligation for comprehensive
disclosure and the ethical obligation for genuine patient understanding.

In the device regulation domain, the comparative analysis confirms that
federal requirements predominate due to federal preemption principles.
However, California's general consumer protection statutes and medical
practice act create supplementary obligations that affect how Physical AI
devices are deployed in clinical settings. The Medical Board of
California's oversight of medical practice extends to situations where
robotic systems perform functions that constitute the practice of medicine,
creating potential regulatory exposure for Physical AI operators who lack
medical licenses. The adapted 21 CFR Part 312 \cite{cfr312-adapt} addresses
this through its Physical AI operator definition, which requires that
robotic procedures be conducted under the supervision of a licensed physician
who maintains clinical responsibility regardless of the degree of robotic
autonomy.

\subsubsection{Robotics-Specific Considerations for Oncology Clearance}

The FDA clearance landscape for surgical robots in oncology presents specific
challenges that Physical AI trial sponsors must navigate. As discussed in
Subsection~\ref{subsec:robotics-regulatory}, robotically-assisted surgical
systems have not received marketing authorization specifically for cancer
prevention or treatment. This means that surgical robots used in oncology
procedures operate under general surgical clearances, with oncology-specific
use constituting off-label deployment in the regulatory sense. Physical AI
oncology trials provide the mechanism for generating the clinical evidence
needed for oncology-specific clearances.

The 510(k) pathway has been the primary route for surgical robot clearance.
The da Vinci Surgical System's original 510(k) clearance established the
predicate for teleoperated surgical robots, and subsequent systems including
Hugo RAS and Versius have been evaluated against this predicate or its
descendants. For Physical AI systems that build on existing teleoperated
platforms by adding AI-assisted features, the 510(k) pathway remains
available if the sponsor can demonstrate substantial equivalence. However,
the determination of substantial equivalence becomes more complex when AI
features change the nature of the surgical interaction. A robot that
autonomously identifies tissue planes, makes real-time surgical decisions,
and adapts its approach based on intraoperative findings operates
fundamentally differently from a passive teleoperation system, even if the
physical hardware is identical.

The De Novo classification pathway offers a potentially appropriate route
for novel Physical AI systems that combine existing robotic hardware with
unprecedented AI autonomy. The De Novo process creates a new device
classification with associated special controls, which could include Physical
AI-specific requirements for AI validation, human oversight, and safety
monitoring. The special controls established through De Novo classification
would become the regulatory standard for subsequent Physical AI devices
seeking 510(k) clearance, effectively creating a new regulatory category
for autonomous surgical robots.

The Breakthrough Device designation program provides expedited development,
assessment, and review for certain devices that offer transformative
advantages for patients with life-threatening or irreversibly debilitating
conditions. Physical AI systems designed for oncology applications may
qualify for Breakthrough designation if they demonstrate the potential to
provide more effective treatment than existing alternatives. The Breakthrough
pathway includes enhanced FDA interaction during development, which could
help sponsors address the novel regulatory questions that Physical AI
presents. The comprehensive evidence base provided by this National Platform,
including simulation data \cite{site-docs}, adapted regulatory standards
\cite{ich-adapt, cfr50-adapt, cfr312-adapt}, and quantitative readiness
scores \cite{usl-standard}, strengthens the case for Breakthrough
designation by demonstrating thorough pre-submission preparation.

The interaction between device clearance pathways and the clinical trial
framework creates a temporal sequencing challenge for Physical AI sponsors.
A Physical AI system that requires De Novo or PMA clearance cannot be
commercially marketed until the clearance process is complete, but the
clinical evidence needed for clearance can only be generated through clinical
trials. The IDE pathway provides the mechanism for conducting these trials
with uncleaned devices, but IDE applications must include sufficient
pre-clinical evidence to support human subjects exposure. The Phase 0
simulation validation requirement in the adapted 21 CFR Part 312
\cite{cfr312-adapt} directly addresses this sequencing challenge by
generating simulation-based safety evidence that can support IDE applications,
enabling clinical trials to proceed while the device clearance process
advances in parallel.

The oncology-specific endpoint requirements add another layer of complexity
to Physical AI device clearance. Standard surgical robot endpoints focus on
procedural success, complication rates, and recovery times. Oncology
endpoints require longer-term follow-up for recurrence, disease-free
survival, and overall survival. This means that Physical AI oncology trials
must be designed with sufficient follow-up duration to generate oncology-
relevant evidence, extending trial timelines beyond what purely procedural
endpoints would require. The adapted ICH E6(R3) protocol requirements
\cite{ich-adapt} address this by mandating that Physical AI trial protocols
specify both procedural endpoints (robot performance, safety metrics) and
oncology endpoints (disease outcomes, survival data), ensuring that the trial
generates evidence sufficient for both device clearance and oncology-specific
marketing authorization.

Table~\ref{tab:clearance-oncology} summarizes the clearance pathway
considerations for each Physical AI system category as classified in the
USL framework \cite{usl-standard} when deployed specifically for oncology
applications.

\begin{table}[H]
\centering
\caption{Physical AI System Categories and Oncology Clearance Considerations}
\label{tab:clearance-oncology}
\small
\begin{tabularx}{\textwidth}{l l X}
\toprule
\textbf{Physical AI Category} & \textbf{Likely Pathway} & \textbf{Oncology-Specific Considerations} \\
\midrule
Surgical robots & 510(k) or De Novo & Existing clearance for general surgery; oncology endpoints require separate evidence \\
Collaborative robots & 510(k) or De Novo & Novel clinical use; specimen handling and lab automation may qualify as non-significant risk \\
Radiotherapy positioning & 510(k) & Existing predicate devices; sub-millimeter accuracy critical for tumor targeting \\
Needle-placement robots & De Novo or PMA & Biopsy targeting for tumor sampling; autonomous guidance raises classification questions \\
Social companion robots & Possible exemption & Low physical risk; patient communication triggers AB 3030 and AB 489 requirements \\
Humanoid robots & De Novo & No predicate exists; risk level depends on permitted patient interactions \\
Radiotherapy tracking & 510(k) & Motion compensation during treatment; real-time adaptation may change classification \\
Imaging robots & 510(k) & Existing imaging device predicates; AI-guided positioning adds complexity \\
Steerable needle robots & De Novo or PMA & Novel technology; autonomous tissue navigation carries significant risk \\
Rehabilitation exoskeletons & 510(k) or De Novo & Existing predicates for rehabilitation; oncology rehabilitation is novel indication \\
\bottomrule
\end{tabularx}
\end{table}

The clearance pathway analysis in Table~\ref{tab:clearance-oncology}
demonstrates that Physical AI oncology trials will likely require multiple
concurrent regulatory pathways, with different Physical AI system categories
following different clearance routes. 

% The National Platform's unified
% approach, which integrates all Physical AI systems under a common regulatory
% framework \cite{ich-adapt, cfr50-adapt, cfr312-adapt} with standardized
% readiness assessment \cite{usl-standard}, reduces the complexity of managing
% multiple concurrent pathways by providing a single documentation and
% compliance infrastructure that supports all device categories simultaneously.

%%% END ADDITIONAL CONTENT FOR SECTION 3 %%%
